\documentclass[11pt]{article}

\usepackage[utf8]{inputenc}
\usepackage[T1]{fontenc}
\usepackage{lmodern}
\usepackage{geometry}
\usepackage{amsmath,amssymb,amsfonts,amsthm,mathtools}
\usepackage{bm}
\usepackage{enumitem}
\usepackage{microtype}
\usepackage{hyperref}
\usepackage{titlesec}
\usepackage{booktabs}
\usepackage{array}
\usepackage{setspace}
\usepackage{mathrsfs}
\usepackage{physics}

\geometry{margin=1in}
\hypersetup{colorlinks=true, linkcolor=black, urlcolor=blue, citecolor=black}
\setlist[enumerate]{topsep=0.4em,itemsep=0.35em}
\setlist[itemize]{topsep=0.3em,itemsep=0.25em}
\titleformat{\section}{\large\bfseries}{\thesection.}{0.5em}{}
\titleformat{\subsection}{\normalsize\bfseries}{\thesubsection.}{0.5em}{}
\setstretch{1.06}

\newcommand{\Aether}{\AE{}ther}
\newcommand{\Aflow}{\AE{}ther-flow}
\newcommand{\TheoryName}{\Aflow{} Interpretation of Relativity}
\newcommand{\TheoryProgram}{\Aether{} / \Aflow{} framework}
\newcommand{\Sub}{\mathcal{A}}
\newcommand{\ReBridge}{g^{\mathrm{br}}}
\newcommand{\OpMetric}{g^{\mathrm{op}}}
\newcommand{\Dperp}{\mathcal D}
\newcommand{\PiTheta}{\Pi^{(\vartheta)}}
\newcommand{\PiQ}{\Pi^{(Q)}}
\newcommand{\PiOmega}{\Pi^{(\Omega)}}
\newcommand{\SigmaOp}{\Sigma^{\mathrm{op}}}
\newcommand{\SigmaLong}{\Sigma^{(\chi,\ge 3)}}
\newcommand{\ReEff}{\widetilde{\mathcal T}_{\mu\nu}^{\mathrm{br,eff}}}
\newcommand{\ReLong}{\widetilde{\mathcal R}_{\mu\nu}^{(\chi,\ge 3)}}
\newcommand{\ReOff}{\widetilde{\mathcal R}_{\mu\nu}^{\mathrm{off}}}
\newcommand{\ReOffTwo}{\widetilde{\mathcal R}_{\mu\nu}^{\mathrm{off},(2)}}
\newcommand{\DeltaTop}{\Delta T_{\mu\nu}^{\mathrm{op}}}
\newcommand{\EinLin}{\mathcal E_{\ReBridge}}
\newcommand{\EinQuad}{\mathcal Q_{\ge 2}^{\mathrm{Ein}}}
\newcommand{\AY}{\mathcal A_Y}
\newcommand{\AZ}{\mathcal A_Z}

\newtheorem{definition}{Definition}
\newtheorem{proposition}{Proposition}
\newtheorem{remark}{Remark}
\newtheorem{corollary}{Corollary}
\newtheorem{theorem}{Theorem}

\title{\textbf{The \TheoryName{}}\\[0.4em]
\large Operational-Metric Redesign Note}
\author{}
\date{}

\begin{document}

\maketitle

\begin{abstract}
The positive quotient reduced-system, theorem, decay, and asymptotic line remains closed and is not reopened here. The two recent redesign notes also remain fixed in status. The bridge-first redesign was only a same-scope longitudinal repair: on the revised bridge branch one still has
\[
G_{\mu\nu}[\ReBridge]
=
8\pi G_N\,T_{\mu\nu}^{\mathrm{matter}}
+\ReEff,
\qquad
\ReEff=\ReOff+\ReLong,
\]
where \(\ReLong=\mathcal O_{\ge 3}\) is the higher-order longitudinal remainder and \(\ReOff\) is the genuine off-longitudinal \(Q/W\)-sector stress. The second matter-route redesign was then also negative because its new correction is matter dependent, vanishes on matter-free reduced data, and therefore cannot remove the same matter-free off-longitudinal remainder.

The present note records the sharpest broader next step. Rather than modifying only the bridge convergence slot or only the matter action, it asks whether the sole operational observer metric itself must absorb the surviving auxiliary stress. The broader redesign is parametrized observer-side by
\[
\OpMetric_{\mu\nu}
=
\ReBridge_{\mu\nu}
+\SigmaOp_{\mu\nu},
\qquad
\SigmaOp_{\mu\nu}
=
\alpha_Q\PiQ_{\mu\nu}
+\alpha_\Omega\PiOmega_{\mu\nu}
+\SigmaLong_{\mu\nu},
\]
where \(\PiQ\) and \(\PiOmega\) are the already-identified off-longitudinal response tensors and \(\SigmaLong=\mathcal O_{\ge 3}\) is any higher-order longitudinal completion term.

Recomputing the eliminated observer equation on this broader redesign gives
\[
G_{\mu\nu}[\OpMetric]
=
8\pi G_N\,T_{\mu\nu}^{\mathrm{matter}}[\OpMetric,\psi]
+\mathcal R_{\mu\nu}^{\mathrm{op}},
\]
with
\[
\mathcal R_{\mu\nu}^{\mathrm{op}}
=
\ReOff+\ReLong
+\EinLin(\SigmaOp)_{\mu\nu}
+\EinQuad[\SigmaOp]_{\mu\nu}
-8\pi G_N\,\DeltaTop.
\]
This identity makes the next proof burden explicit. On matter-free reduced data, \(\DeltaTop=0\), so the redesign acts only if the operational-metric shift itself contributes. The two exact targets are therefore:
\begin{enumerate}
    \item remove the off-longitudinal remainder by forcing
    \[
    \EinLin(\alpha_Q\PiQ+\alpha_\Omega\PiOmega)_{\mu\nu}
    =
    -\ReOffTwo
    \]
    on matter-free reduced data;
    \item remove, or at least structurally explain, the higher-order longitudinal remainder by forcing the remaining longitudinal sector to vanish or to be only gauge/constraint exact after the Einstein constraints and slice equations are imposed.
\end{enumerate}
The note does not claim that such a redesign has already been found. Its narrower result is to show why the next honest branch must be broader than the two failed minimal repairs and to reduce that broader redesign to an explicit operational-metric tensor target.
\end{abstract}

\section{Introduction}

The active derivational line is now sharp about what is and is not still open. The explicit ordered-flow-label quotient candidate, the quotient-architecture screen, the quotient reduced-system note, the quotient local/coercive note, the quotient theorem note, and the quotient decay and asymptotic notes already closed the positive quotient dynamical line at the recorded reduced-system, theorem, decay, and asymptotic scopes. None of that work is reopened here.

The first bridge-first redesign and the second matter-route redesign also now have disciplined status. The convergence-response bridge change removed only the old same-scope longitudinal slot and left both a higher-order longitudinal remainder and a genuine off-longitudinal observer remainder. The smallest honest matter-route-only change then remained negative because its correction is matter dependent and therefore vanishes on matter-free reduced data while the off-longitudinal remainder survives anyway.

\paragraph{Framework and claim boundary.}
\TheoryProgram{} denotes the broader \Aether{} / \Aflow{} framework built on the claims that the \Aether{} is the underlying four-dimensional substrate of reality, the \Aflow{} is its intrinsic ordered motion, observed three-dimensional space is the local experiential slice of that deeper substrate, S-time is the experienced order of change arising from matter, light, and the \Aflow{}, and observed expansion is the three-dimensional appearance of deeper four-dimensional ordered motion. Within that framework, \TheoryName{} denotes the exact-closure theory in which the effective gravitational dynamics are adopted to be exactly Einsteinian with universal matter coupling. In the active sequence, \emph{adoption} means use of that established relativistic dynamics without claiming substrate derivation, while \emph{derivation} is reserved for a first-principles recovery from explicit substrate variables.

The present note stays strictly inside that boundary. It does not claim a completed substrate derivation. It does not reopen the positive quotient PDE line. It treats the two recent redesign notes as failed minimal repairs, defines the broader operational-metric or bridge-map redesign class at observer level, recomputes the eliminated observer equation in that broader class, and fixes the exact targets that any honest next candidate must hit.

\section{What Is Closed and What the Minimal Repairs Failed To Do}

\begin{definition}[Current redesigned bridge branch]
\label{def:current}
The \emph{current redesigned bridge branch} is the branch already fixed by the redesigned-bridge closure/tensor-verdict note:
\begin{enumerate}
    \item the quotient equations
    \begin{equation}
    \AY=0,
    \qquad
    \AZ=0,
    \label{eq:AYAZ}
    \end{equation}
    are imposed from the outset;
    \item the bridge-first redesign has already been absorbed into the current operational bridge metric \(\ReBridge_{\mu\nu}\);
    \item matter is still coupled through that same single operational metric;
    \item the full eliminated observer equation already recorded on that branch is
    \begin{equation}
    G_{\mu\nu}[\ReBridge]
    =
    8\pi G_N\,T_{\mu\nu}^{\mathrm{matter}}
    +\ReEff,
    \qquad
    \ReEff=\ReOff+\ReLong.
    \label{eq:oldobs}
    \end{equation}
\end{enumerate}
\end{definition}

\begin{remark}
Definition \ref{def:current} keeps the positive quotient reduced-system, theorem, decay, and asymptotic line closed. The only live burden is the surviving observer remainder in \eqref{eq:oldobs}.
\end{remark}

\begin{proposition}[Bridge-first redesign was only a same-scope longitudinal repair]
\label{prop:bridgefail}
On the current redesigned bridge branch:
\begin{enumerate}
    \item the longitudinal contribution satisfies
    \begin{equation}
    \ReLong=\mathcal O_{\ge 3},
    \label{eq:long}
    \end{equation}
    so the bridge-first redesign does not supply an all-order longitudinal identity;
    \item the off-longitudinal remainder \(\ReOff\) survives genuinely;
    \item in particular, on the rotationally nontrivial convergence-free matter-free regime
    \begin{equation}
    K\equiv 0,
    \qquad
    \Omega_{ij}^T\not\equiv 0,
    \qquad
    T_{\mu\nu}^{\mathrm{matter}}=0,
    \label{eq:rot}
    \end{equation}
    one still has
    \begin{equation}
    G_{\mu\nu}[\ReBridge]
    =
    T_{\mu\nu}^{(W,\mathrm{br,eff},2)}
    +\mathcal O_{\ge 3},
    \qquad
    \rho_W^{\mathrm{br},(2)}
    =
    \frac{\kappa_\Omega}{4}\abs{\Omega^T}^2.
    \label{eq:oldW}
    \end{equation}
\end{enumerate}
\end{proposition}

\begin{proof}
This is exactly the content of the redesigned-bridge closure/tensor-verdict note. The convergence-response bridge term cancels only the old quadratic longitudinal slot and is blind, at the same scope, to the decisive off-longitudinal \(Q/W\)-sector remainder.
\end{proof}

\begin{proposition}[Matter-route-only redesign is also negative]
\label{prop:matterfail}
The second matter-route redesign does not remove the surviving observer remainder. Its new correction is matter dependent, so on the matter-free regime \eqref{eq:rot} it vanishes and the observer equation reduces again to \eqref{eq:oldW}.
\end{proposition}

\begin{proof}
This is the theorem-level verdict of the second matter-route redesign note. The modified matter route changes the observer equation only through a dressed matter stress correction proportional to the matter sector itself. Hence it is inactive on matter-free reduced data.
\end{proof}

\begin{corollary}[The next redesign must act on matter-free reduced data]
\label{cor:matterfree}
Any honest next redesign must change the observer equation even when \(T_{\mu\nu}^{\mathrm{matter}}=0\). In particular, it cannot be a matter-route-only dressing whose effect disappears on the matter-free regime \eqref{eq:rot}.
\end{corollary}

\begin{proof}
Combine Propositions \ref{prop:bridgefail} and \ref{prop:matterfail}.
\end{proof}

\section{Broader Operational-Metric or Bridge-Map Redesign}

The previous note already identified the key logical fork. If one promotes the dressed matter metric to the sole operational metric for clocks, light, the Einstein--Hilbert variation, and matter, then one is no longer making a matter-route-only modification. One is changing the operational observer metric itself. That is the broader redesign class opened here.

\begin{definition}[Observer-side response tensors]
\label{def:responses}
On the current redesigned bridge branch, define the already-identified off-longitudinal response tensors by
\begin{equation}
\PiQ_{\mu\nu}
\equiv
D_\mu Q^I\,D_\nu Q_I
-\frac12 \ReBridge_{\mu\nu}
\left(
D_\alpha Q^I\,D^\alpha Q_I+\widehat M_{IJ}Q^IQ^J
\right),
\label{eq:piQ}
\end{equation}
\begin{equation}
\PiOmega_{\mu\nu}
\equiv
\Omega_{\mu\alpha}^T\Omega_\nu^{T\,\alpha}
-\frac14 \ReBridge_{\mu\nu}\Omega_{\alpha\beta}^T\Omega^{T,\alpha\beta}.
\label{eq:piOmega}
\end{equation}
Let \(\SigmaLong_{\mu\nu}\) denote any covariant higher-order longitudinal completion tensor built from the same reduced variables and satisfying
\begin{equation}
\SigmaLong_{\mu\nu}=\mathcal O_{\ge 3}.
\label{eq:sigmalong}
\end{equation}
\end{definition}

\begin{remark}
The tensors in Definition \ref{def:responses} are not new propagating sectors. They are the observer-side stress channels already isolated by the redesigned-bridge and second matter-route verdicts. The point of the broader redesign is to let the operational metric itself see them.
\end{remark}

\begin{definition}[Broader operational-metric redesign]
\label{def:opmetric}
The \emph{broader operational-metric redesign} is the observer-level ansatz
\begin{equation}
\OpMetric_{\mu\nu}
=
\ReBridge_{\mu\nu}
+\SigmaOp_{\mu\nu},
\label{eq:opmetric}
\end{equation}
with
\begin{equation}
\SigmaOp_{\mu\nu}
=
\alpha_Q\PiQ_{\mu\nu}
+\alpha_\Omega\PiOmega_{\mu\nu}
+\SigmaLong_{\mu\nu}.
\label{eq:sigmaop}
\end{equation}
The Einstein--Hilbert term, causal structure, proper-time rule, light propagation, and matter action are all taken to use the same metric \(\OpMetric_{\mu\nu}\).
\end{definition}

\begin{remark}
Definition \ref{def:opmetric} is broader than both recent failed repairs. Unlike the bridge-first redesign, it is not restricted to the convergence channel \(\PiTheta\). Unlike the matter-route-only redesign, it remains active on matter-free reduced data because the metric entering the Einstein tensor has itself changed.
\end{remark}

\begin{proposition}[Observer-side parametrization is enough for the next test]
\label{prop:param}
Any broader bridge-map redesign that still yields one sole operational observer metric can be tested, at the present manuscript scope, through the induced observer-side shift \(\SigmaOp_{\mu\nu}\) of Definition \ref{def:opmetric}.
\end{proposition}

\begin{proof}
The observer equation depends on the broader redesign only through the operational metric used by the Einstein--Hilbert and matter sectors. Whatever bridge-map or pullback change is made substrate-side, its next observer-level test is whether the induced metric correction \(\SigmaOp_{\mu\nu}\) removes the recorded remainder \(\ReEff\). That is exactly the scope of the present note.
\end{proof}

\section{Recomputed Eliminated Observer Equation}

\begin{definition}[Linearized Einstein response around the current redesigned bridge]
\label{def:einlin}
For any symmetric tensor \(H_{\mu\nu}\), define the linearized Einstein response at \(\ReBridge\) by
\begin{equation}
\EinLin(H)_{\mu\nu}
\equiv
\left.
\frac{d}{d\varepsilon}
G_{\mu\nu}[\ReBridge+\varepsilon H]
\right|_{\varepsilon=0}.
\label{eq:einlin}
\end{equation}
Define the nonlinear remainder by
\begin{equation}
\EinQuad[H]_{\mu\nu}
\equiv
G_{\mu\nu}[\ReBridge+H]
-G_{\mu\nu}[\ReBridge]
-\EinLin(H)_{\mu\nu}.
\label{eq:einquad}
\end{equation}
\end{definition}

\begin{definition}[Matter stress correction on the broader redesign]
\label{def:deltat}
Let
\begin{equation}
\DeltaTop
\equiv
T_{\mu\nu}^{\mathrm{matter}}[\OpMetric,\psi]
-T_{\mu\nu}^{\mathrm{matter}}[\ReBridge,\psi].
\label{eq:deltat}
\end{equation}
\end{definition}

\begin{remark}
On matter-free reduced data, \(T_{\mu\nu}^{\mathrm{matter}}=0\) implies \(\DeltaTop=0\). This is the structural difference from the failed matter-route-only redesign.
\end{remark}

\begin{proposition}[Broader redesign observer equation]
\label{prop:observer}
On reduced solutions of the current redesigned bridge branch, the broader operational-metric redesign of Definition \ref{def:opmetric} satisfies
\begin{equation}
G_{\mu\nu}[\OpMetric]
=
8\pi G_N\,T_{\mu\nu}^{\mathrm{matter}}[\OpMetric,\psi]
+\mathcal R_{\mu\nu}^{\mathrm{op}},
\label{eq:observer}
\end{equation}
where
\begin{equation}
\mathcal R_{\mu\nu}^{\mathrm{op}}
=
\ReOff+\ReLong
+\EinLin(\SigmaOp)_{\mu\nu}
+\EinQuad[\SigmaOp]_{\mu\nu}
-8\pi G_N\,\DeltaTop.
\label{eq:observer2}
\end{equation}
\end{proposition}

\begin{proof}
By Definition \ref{def:current},
\[
G_{\mu\nu}[\ReBridge]
=
8\pi G_N\,T_{\mu\nu}^{\mathrm{matter}}[\ReBridge,\psi]
+\ReOff+\ReLong.
\]
By Definition \ref{def:einlin},
\[
G_{\mu\nu}[\OpMetric]
=
G_{\mu\nu}[\ReBridge]
+\EinLin(\SigmaOp)_{\mu\nu}
+\EinQuad[\SigmaOp]_{\mu\nu}.
\]
By Definition \ref{def:deltat},
\[
T_{\mu\nu}^{\mathrm{matter}}[\ReBridge,\psi]
=
T_{\mu\nu}^{\mathrm{matter}}[\OpMetric,\psi]
-\DeltaTop.
\]
Substituting these identities gives \eqref{eq:observer}--\eqref{eq:observer2}.
\end{proof}

\begin{corollary}[Why the broader redesign is the first honest next branch]
\label{cor:honest}
On matter-free reduced data,
\begin{equation}
\mathcal R_{\mu\nu}^{\mathrm{op}}
=
\ReOff+\ReLong
+\EinLin(\SigmaOp)_{\mu\nu}
+\EinQuad[\SigmaOp]_{\mu\nu},
\label{eq:matterfreeobs}
\end{equation}
so the broader redesign acts precisely where the matter-route-only redesign was inactive.
\end{corollary}

\begin{proof}
Set \(T_{\mu\nu}^{\mathrm{matter}}=0\) in Definition \ref{def:deltat} and use Proposition \ref{prop:observer}.
\end{proof}

\section{The Two Explicit Targets}

\begin{definition}[Off-longitudinal cancellation target]
\label{def:offtarget}
The \emph{off-longitudinal cancellation target} is the requirement that, on matter-free reduced data and at the first nontrivial eliminated order,
\begin{equation}
\EinLin(\alpha_Q\PiQ+\alpha_\Omega\PiOmega)_{\mu\nu}
=
-\ReOffTwo.
\label{eq:offtarget}
\end{equation}
\end{definition}

\begin{remark}
Equation \eqref{eq:offtarget} is the first decisive test because the rotationally nontrivial convergence-free regime already shows that the surviving obstruction is genuinely matter free and off longitudinal.
\end{remark}

\begin{definition}[Higher-order longitudinal removal or explanation target]
\label{def:longtarget}
The \emph{higher-order longitudinal target} is the requirement that the remaining longitudinal sector
\begin{equation}
\mathcal L_{\mu\nu}^{\mathrm{long}}
\equiv
\ReLong
+\EinLin(\SigmaLong)_{\mu\nu}
+\EinQuad[\SigmaOp]_{\mu\nu}\Big|_{\mathrm{long}}
\label{eq:longtarget}
\end{equation}
either:
\begin{enumerate}
    \item vanishes identically, or
    \item is only gauge/constraint exact after the Einstein constraints and the reduced slice equations are imposed.
\end{enumerate}
\end{definition}

\begin{definition}[Gauge/constraint exactness at the present scope]
\label{def:pure}
A symmetric tensor \(R_{\mu\nu}\) is \emph{gauge/constraint exact at the present scope} if there exist a one-form \(Y_\mu\), the Einstein Hamiltonian and momentum constraints \(\mathcal H\), \(\mathcal M_\mu\), and reduced slice equations \(E_{Q,I}=0\), \(E_\chi=0\), \(E_{W,i}=0\) such that
\begin{equation}
R_{\mu\nu}
=
\nabla_{(\mu}Y_{\nu)}
+\mathcal P_{\mu\nu}\mathcal H
+\mathcal P_{\mu\nu}{}^\rho\mathcal M_\rho
+\mathcal P_{\mu\nu}^{(Q)I}E_{Q,I}
+\mathcal P_{\mu\nu}^{(\chi)}E_\chi
+\mathcal P_{\mu\nu}^{(W)i}E_{W,i},
\label{eq:pure}
\end{equation}
for differential operators \(\mathcal P\) built from the reduced variables.
\end{definition}

\begin{theorem}[Operational-metric redesign criterion]
\label{thm:criterion}
The broader operational-metric redesign of Definition \ref{def:opmetric} closes the observer equation on the current redesigned bridge branch only if both of the following hold:
\begin{enumerate}
    \item the off-longitudinal cancellation target \eqref{eq:offtarget} is satisfied on matter-free reduced data;
    \item the higher-order longitudinal target of Definition \ref{def:longtarget} is satisfied.
\end{enumerate}
Conversely, if either target fails, then the broader redesign is still negative.
\end{theorem}

\begin{proof}
By Corollary \ref{cor:honest}, the matter-free observer equation on the broader redesign is
\[
G_{\mu\nu}[\OpMetric]
=
\ReOff+\ReLong
+\EinLin(\SigmaOp)_{\mu\nu}
+\EinQuad[\SigmaOp]_{\mu\nu}.
\]
Since \(\SigmaLong=\mathcal O_{\ge 3}\), the only terms that can cancel the quadratic off-longitudinal remainder \(\ReOffTwo\) are the \(\PiQ\)- and \(\PiOmega\)-parts of \(\SigmaOp\) through the linearized Einstein response. This gives target (1). Once that decisive off-longitudinal piece is addressed, the remaining longitudinal sector is exactly \(\mathcal L_{\mu\nu}^{\mathrm{long}}\) of \eqref{eq:longtarget}; exact observer closure requires that sector either to vanish or to be pure gauge/constraint exact in the sense of Definition \ref{def:pure}. If either requirement fails, the residual tensor \(\mathcal R_{\mu\nu}^{\mathrm{op}}\) remains nonzero in the observer equation, so the broader redesign is negative.
\end{proof}

\begin{corollary}[The next manuscript burden after this note]
\label{cor:next}
The next honest manuscript burden is no longer another quotient PDE note, another bridge-only convergence-slot note, or another matter-route-only note. It is an explicit operational-metric or equivalent bridge-map test that computes whether some concrete \(\SigmaOp_{\mu\nu}\) satisfies the two targets of Theorem \ref{thm:criterion}.
\end{corollary}

\begin{proof}
The positive quotient reduced-system, theorem, decay, and asymptotic line is already closed, and Propositions \ref{prop:bridgefail} and \ref{prop:matterfail} classify the two recent minimal repairs as negative. Theorem \ref{thm:criterion} then fixes the only remaining honest proof burden.
\end{proof}

\section{Conclusion}

The active sequence is now sharper again. The positive quotient reduced-system, theorem, decay, and asymptotic line stays closed. The bridge-first redesign and the matter-route-only redesign are both still worth having, but only as failed minimal repairs: the first removed only the old same-scope longitudinal slot, and the second remained matter dependent and therefore inactive on the decisive matter-free data.

The broader next step is now explicit. If the exact-closure goal is to remain active, the redesign must act on matter-free reduced data and must see the off-longitudinal \(Q/W\) sectors as well as the higher-order longitudinal remainder. At the present manuscript scope, that burden is cleanly expressed by asking whether the sole operational observer metric itself must absorb the surviving auxiliary stress. The recomputed observer equation \eqref{eq:observer} shows exactly what that means. The next candidate must supply an operational-metric shift whose Einstein response cancels the matter-free off-longitudinal remainder and removes, or at least structurally explains, the higher-order longitudinal remainder. That is now the sharpest honest next derivational test.

\end{document}
